\scanpage{131}
\mspage{64}

We find that the composites
\[
T_\pi \xrightarrow{\sim} L_i \longrightarrow
H^2(\pi,\widehat{\mathbb Z})^{\otimes-1}
\]
do not depend on the choice of $i$, so that the orientation module $T_\pi$ of the group \inserted{profinite $\pi$} with loops $\pi$ is identified with the \uline{dual} of $H^2(\pi,\widehat{\mathbb Z})$ (if $g\ne0$).

Let us now consider a homomorphism of groups with loops
\[
\begin{gathered}
\pi'\xrightarrow{f}\pi\\
I'_0\xrightarrow{\tau}I
\end{gathered}
\]
associated with an isomorphism $T_{\pi'}\simeq T_\pi$ and a degree map
$d:I'_0\longrightarrow\mathbb N$
\deleted{(\unsure{integers} \illegible{3}}
(N.B. $i'\in I'\setminus I'_0\longrightarrow f(L'_{i'})=(1)$.

We would like to deduce from this a diagram of morphisms of topoi
\diagram{\begin{tikzcd}
B_{\Pi_1 U'} \arrow[r,dash] \arrow[d] & B_{\deleted{\unsure{\Pi_1}}X',U'} \arrow[r,dash] \arrow[d] & B_{\Pi_1 X'} \arrow[d] \\
B_{\Pi_1 U} \arrow[r,dash] & B_{\deleted{\unsure{\Pi_1}}X,U} \arrow[r,dash] & B_{\Pi_1 X}
\end{tikzcd}}
and a homomorphism \uline{trace} on the cohomology with proper supports of $U',U$ (defined in terms of cohomology on $B_{X',U'}, B_{X,U}$ relatively)
\deleted{\unsure{which an isomorphism}} \unsure{inducing} \deleted{\unsure{for}} an \uline{isomorphism}
\[
H^2_!(U',\widehat{\mathbb Z})\longrightarrow
H^2_!(U,\widehat{\mathbb Z})
\]

\scanpage{132}

which is precisely \inserted{(\unsure{contragredient|contragredient} of)} (the isomorphism of orientation modules associated with $f$....

14. Where we return to \uline{mixed} morphisms (corresponding, in the topological setting, to the case of
$f:(X',S')\longrightarrow(X,S)$ such that we have
$[f(X'\setminus S')\subset X\setminus S$, i.e.\ \inserted{$S'\supset f^{-1}(S)$,} \inserted{but}$]$
\deleted{\illegible{1}} ($S'\ne f^{-1}(S)$, i.e.\ there are points of $\unsure{S'|S}$ that are sent into $U=X\setminus S$.
\margin{$X,X'$ $T$-oriented}

In the topossic and groupoidal \deleted{language} \inserted{paradigm}, $(X,S)$ is \deleted{defined} \inserted{described} by a diagram of groupoids
\[
\Pi_{D^*}\longrightarrow\Pi_U
\qquad
(\text{and }k:T_{\Pi_{D^*}}\simeq
\text{ local system of the }\pi_1\text{ on }\Pi_{D^*})
\]
\deleted{\illegible{1}}
or of topoi
\[
B_{D^*}\longrightarrow B_U
\]
and \deleted{\unsure{an}} isomorphism \inserted{$k$} of the local system of the $\pi_1$ on $B_{D^*}$, into the \inserted{locally} constant system $T$,
\deleted{\illegible{4} $(\deleted{\unsure{X}},U')$}
\inserted{which allows us to define $B_D\;(=B_{\unsure{\pi_0 D_*}})$ (the discrete topos)}
and the topos $B_{X,U}$, fitting into the diagram \deleted{$\unsure{B_{\Pi_1 D^*}}$} of topoi

\scanpage{133}
\mspage{65}

\diagram{\begin{tikzcd}[column sep=large]
& B_{D^*}\arrow[dl,"\rho"']\arrow[dr,"\sigma"] & \\
B_U\arrow[dr,"\varphi"'] && B_D\; (\text{discrete topos }\xleftarrow{\sim} B_S)\arrow[dl,"\psi"] \\
& B_{X,U} &
\end{tikzcd}}
where $\varphi,\psi$ are morphisms of \unsure{inclusion|induction} of subtopoi, $\varphi$ open $\psi$ closed, complementary to one another, and likewise for $(X',S')$, described by an analogous diagram
\diagram{\begin{tikzcd}
& B_{D'^*}\arrow[dl]\arrow[dr] & \\
B_{U'}\arrow[dr] && B_{D'}\simeq\unsure{B_{S'}|B_S}\arrow[dl] \\
& B_{X',U'} &
\end{tikzcd}}

That said, we absolutely want, in a description of
$f:(X',S')\longrightarrow(X,S)$, for it to allow us to recover not only $B_{U'}\longrightarrow B_U$ (which will be secured by giving an $f_\pi:\pi'\longrightarrow\pi$), but also
\[
B_{X',U'}\longrightarrow B_{X,U}\qquad\unsure{g|q}.
\]
We will have \deleted{\unsure{$S$}}
\[
S'=\underset{\substack{\Vert\\f^{-1}(S)}}{S'_0}\amalg S'_1,
\quad\text{hence}\quad D'=D'_0\amalg D'_1,
\]

\scanpage{134}
\mspage{66}

and in fact $f$ induces $S'_1\xrightarrow{f}U$ (and not $S'_1\longrightarrow S$) which must be made explicit, at the level \deleted{\illegible{2}} of multi-Galois topoi, by a morphism
\diagram{\begin{tikzcd}
B_{D'_1}\arrow[r] & B_U \\
B_{S'_1}\arrow[u,"\simeq",draw=none] &
\end{tikzcd}}
i.e.\ by giving a family of universal coverings of $U$, parameterized by $S'_1=I'_1$, or (\unsure{if such a} universal covering is chosen, $U$ being connected, hence a $\pi=\operatorname{Aut}(\widetilde U)$) by a family of torsors under \deleted{\unsure{the}} $\pi$, $(P_{i'})_{i'\in I'_1}$.
This being \deleted{\unsure{defined}} \inserted{laid down} we \unsure{will be able to} describe, in terms of
\diagram{\begin{tikzcd}[ampersand replacement=\&,column sep=3.5em,row sep=3em]
B_{D_0'{}^*} \arrow[r,no head] \arrow[d] \& B_{U'} \arrow[d] \\
B_{D^*} \arrow[r,no head]
\arrow[ur,bend left=30,shorten <=1.5em,shorten >=1.5em,"\sim"'] \& B_U
\end{tikzcd}}
and of
\[
B_{D_1'^*}\longrightarrow B_{D'_1}\longrightarrow B_U
\]
\deleted{\unsure{a}} the morphism of topoi $B_{X',U'}\longrightarrow B_{X,U}$
% a small double tick stands before $B_{X',U'}$ (below $B_{D'_1}$ of the display above): probably the call of the marginal note
--- I pass over the details of the description
\margin{It seems that we are entitled to give the description of the morphisms of topoi $B_{X',U'}\longrightarrow B_{X,U}$ that induce $U'\longrightarrow U$ and that on $D_1'^*$ \unsure{$(=D'^*\vert U$} \deleted{\illegible{2}}, factor through $D_1'^*\longrightarrow D'_1$...}

\scanpage{135}

When we give ourselves a group of operators \deleted{$\Sigma$}$\Gamma$ on the situation $(X',S')\longrightarrow(X,S)$ hence on the groupoidal or topossic situation, we must take this into account in the description above.
Thus $I'_0$ and $I'_1$ are stable under $\Gamma$, and the morphism \deleted{\illegible{1}} $B_{\unsure{D'_1}}\longrightarrow B_U$ must be stable under $\Gamma$.
\deleted{So we have a \unsure{nice} set $\unsure{E_f=\bigcup E_i}$ ---}
\deleted{In terms of a universal covering $\widetilde U$ of $U$, let us consider the \deleted{\illegible{1}} sets $\unsure{\operatorname{Hom}}(f_*(\unsure{s'_{i'}}),\widetilde U)$}
Which also means, probably, that we have a morphism of topoi from
$B_{D'_1,\Sigma}=B_{I'_1,\unsure{\Sigma|\Gamma}}$
into $B_{U,\deleted{\Sigma}\Gamma}$.
Choosing for every orbit of $\Gamma$ in $I'_1$ a representative $s'\in I'_1$, and considering its stabilizer $\Gamma_{s'}\subset\Gamma$, we must therefore for every such orbit i.e.\ \deleted{\unsure{every}} every $i'$ give ourselves a $\Gamma_{s'}$-object of the category of universal coverings of $(U,\Gamma)$.
\unsure{Simultaneously|Equivalently} \deleted{therefore} Describing $(U,\Gamma)$ in terms

\scanpage{136}
\mspage{67}

of an extension $\mathcal E$ of $\Sigma\subset\Gamma$ by $\pi$ (by choosing a universal covering $\widetilde U$ of $U$), the datum of $B_{\Gamma_{s'}}\longrightarrow B_{X,U}$ compatible with everything amounts unless I am mistaken to the datum of a right torsor $P$ under $\pi$ (allowing us to twist the \unsure{reference} universal covering $\widetilde U$ of $\unsure{U^\Gamma|U,}$ with $P$ \deleted{and} \inserted{hence also to twist $\mathcal E$ by $P$}), and of a splitting of $\mathcal E^P\longrightarrow\Gamma$ over $\Gamma_{s'}$.
[The isomorphism class of such an object $\longrightarrow$ a $\pi$-conjugacy class of splittings of $\mathcal E\longrightarrow\Gamma$ over $\Gamma_{s'}$... ).

An automorphism of such a pair \deleted{corresponds} $P,\ \Gamma_{s'}\xrightarrow{q}\mathcal E^P$ corresponds to an $\alpha\in\pi$ that \uline{centralizes} $\Gamma_{s'}$, \inserted{i.e.\ that is fixed by $\Gamma_{s'}$ operating on $\pi$.} (If the $\Gamma_{s'}$ operate strongly enough on $\pi$ for us to know that $\pi^{\Gamma_{s'}}=(1)$, the object $(P,q:\Gamma_{s'}\longrightarrow\mathcal E^P)$ is defined up to \uline{unique isomorphism} by the isomorphism class \deleted{\unsure{of it}} (class of splittings of $\mathcal E$ over $\Gamma_{s'}$).
\margin{\unsure{!!!|[three vertical strokes]}}
It seems probable to me that this is always the case in the \uline{arithmetic} case, where $\Gamma_{s'}$ is an open subgroup of the Galois group $\Gamma_{\unsure{k|K}}$ (whose outer operation is \inserted{draconian !}).

\scanpage{137}

This goes in the direction of the conjecture that a splitting class of $\mathcal E$ over $\Gamma$ is ``as good'' as a rational point of $U$ over $K$ --- and allows us to paradigm what such a rational point would allow us to obtain....

\textcircled{15} \uline{Return to the topological case} Let $\Gamma$ be a finite group operating on a group with loops $\pi$ --- let us suppose that this comes from a topological situation, $\Gamma$ operating on $(X,S)$. There are points to check (\uline{\unsure{anabelian|anabelian}} cases)
\margin{Structure of the critical $\Gamma$-orbits in terms of extension...}

\[
\text{a) Operation of }\Gamma\text{ trivial}
\quad\Longleftrightarrow\quad
\underset{\text{trivial}}{\Gamma\longrightarrow\operatorname{Autext}(\pi)}
\]
in other words : an automorphism of finite order of $(X,S)$ can be isotopic to the identity (or even \inserted{merely} homotopic, it comes to the same thing anyway) only if it is trivial.

This is even known in ``abstract'' Algebraic Geometry, once we admit that $u$ \unsure{preserves|preserves} a \deleted{\unsure{complex}} \inserted{complex} structure, \deleted{\unsure{conformal}} \inserted{\unsure{this is precisely}}

\scanpage{138}
\mspage{68}

so if we \deleted{suppose} \inserted{admit} that there is no wildness.... Probably every action of a finite group $\Gamma$ leaves conformal structure invariant, \deleted{\illegible{1}} hence its restriction to
\[
\Gamma^0=\operatorname{Ker}(\Gamma\xrightarrow{\chi}\{\pm1\})
\]
leaves complex structure invariant.

Let us suppose from now on that $\Gamma$ operates faithfully, \inserted{$\Gamma\ne1$,} \deleted{\illegible{1}} The choice of a complex structure on $Y=X/\Gamma^0$ \inserted{defines one on $X$;} (N.B.: \deleted{\illegible{1}} $\Gamma/\Gamma^0$ still operates \inserted{on $X/\Gamma^0$} in fact \deleted{$X/\Gamma^0$} \unsure{if} $\Gamma\ne\Gamma^0$ i.e.\ $\Gamma/\Gamma^0\simeq\{\pm1\}$, $X/\Gamma^0$ is endowed with a ``structure of a \deleted{\illegible{2}} real algebraic curve \inserted{conformal structure}''... ) stable under $\Gamma^0$ --- and (stable under $\Gamma$ if we choose that of $Y$ \deleted{\unsure{invariant}} \inserted{left invariant} under \unsure{the element|the element} nontrivial in $\Gamma/\Gamma^0$ (if there is one).).
\margin{\unsure{works|as} every time we have a ramified covering of conformal surfaces}

\deleted{Now} Every $x\in U^\Gamma$ defines a class \deleted{\illegible{1}} \inserted{\deleted{\unsure{class}} of $\pi$-conjugacy of} splittings of the extension $\mathcal E$ of $\Gamma$ by $\pi$, as we see by taking $x$ as base point. ! Note that $U^\Gamma$ is a finite set \inserted{if $\Gamma^0\ne1$, [i.e.\ $\Gamma$ not reduced to an \deleted{\illegible{1}} \unsure{anti-involution} \illegible{1})}.
\margin{\uline{N.B.} $U^\Gamma\ne\varnothing$ implies that $\Gamma^0$ is cyclic...}
We find \deleted{therefore} a map
\[
U^\Gamma\longrightarrow\text{set of classes of }\pi\text{-conjugacy of splittings of }\mathcal E\longrightarrow\Gamma,
\]

\scanpage{139}
\mspage{69}

i.e.\ \unsure{set|to the} of liftings of $\Gamma\longrightarrow\operatorname{Autext}\pi$ \inserted{to} $\Gamma\longrightarrow\operatorname{Aut}(\pi)$, modulo $\pi$-conjugation.

\uline{Question} \inserted{(If $\Gamma^0\ne(1)$)} Is this map bijective ??

[If $\Gamma^0=\{1\}$, $\Gamma/\Gamma^0\simeq\{\pm1\}$, then $X^\Gamma$ is the set of real points of a real algebraic curve and $U^\Gamma=X^\Gamma-S^\Gamma$ is the complement of a finite subset in it, we have
\[
\pi_0(X^\Gamma-S^\Gamma)\longrightarrow
\pi\text{-classes of splittings of }\mathcal E\text{ over }\Gamma
\]
and the \deleted{\illegible{1}} \uline{analogous} question of bijectivity arises, for extensions of $\pm1$ by $\pi$... ]

For the injectivity of the map in the case $\Gamma^0\ne(1)$, we may suppose $\Gamma=\unsure{\Gamma^0_a|\Gamma^0}\simeq\mathbb Z/p\mathbb Z$, with $p$ prime, i.e.\ $\Gamma$ generated by a complex automorphism \inserted{$\unsure{\gamma|y}$} of order $p$, which defines (if $x,y\in U^\Gamma$, $x\ne y$) an automorphism of order $p$ of $\pi_1(U,x)$, \deleted{\unsure{say}} $u_x$, \unsure{that is} \deleted{a} (hence $\pi$-conjugacy class of automorphisms of $\pi$ of order $p$), and another $u_y$.
We have to prove that $u_x,u_y$ are not conjugate under $\pi$

\scanpage{140}

\uline{Theorem} Let $\unsure{U|X}$ be \deleted{\unsure{a surface}} a connected topological \inserted{space} \inserted{(\deleted{\illegible{2}}, by paths)}, $\Gamma$ a finite group operating \deleted{faithfully} on $U$, $\widetilde U$ a universal covering of $U$, hence \deleted{operating} an extension group
\[
1\longrightarrow\pi\xrightarrow{i}\mathcal E\xrightarrow{p}\Gamma\longrightarrow1
\]
operating faithfully on $\widetilde U$ ($\pi=\operatorname{Aut}_U(\widetilde U)\simeq\pi_1(U)$... ).
For every fixed point $x\in U^\Gamma$, $\Gamma$ operates on the universal covering \unsure{punctured|pointed} over $x$, say \deleted{\illegible{1}} $\unsure{R_x|\mathcal R_x}$, leaving fixed the point \inserted{marked} $\widetilde x$ of $\unsure{R_x}$ over $x$, hence a splitting of the extension relative to $\unsure{R_x}$
\diagram{\begin{tikzcd}
1\arrow[r] & \pi_x\arrow[r,"i_x"] & \mathcal E_x\arrow[r,"p_x"] & \Gamma\arrow[r]\arrow[l,bend left=40,"\sigma_x"] & 1
\end{tikzcd}}
and for every \deleted{choice of path} \inserted{isomorphism} $\ell$ \inserted{(``path'')} $\unsure{R_x}\simeq\widetilde U$, inducing an isomorphism of extensions $\mathcal E_x\xrightarrow{\sim}\mathcal E$ (defined \unsure{of|in} modulo composition with an inner automorphism induced by an $\alpha\in\pi$... ) we find \uline{by transport} of structure a splitting $\sigma_{x,\ell}$ of $\mathcal E\xrightarrow{p}\Gamma$, (which, for varying $\ell$, is defined up to an inner automorphism by an $\alpha\in\pi$).
We thus find a map
\[
(*)\quad U^\Gamma\longrightarrow
\begin{gathered}
\text{classes of }\pi\text{-conjugacy of splittings}\\
\text{of the extension }\mathcal E\text{ of }\Gamma\text{ by }\pi\\
=\text{classes of }\pi\text{-conjugacy of subgroups }\inserted{\unsure{sections}}\ \Gamma'\text{ of }\mathcal E
\end{gathered}
\]
